\section{Ergebnisse}\label{chap:ergebnisse}

Im Folgenden werden die Benchmark-Ergebnisse präsentiert. Für das erste Szenario wird ein
vollständiger Ergebnissatz aufgeführt (\textit{Mark}, Laufzeit, Hauptspeicher und VRAM). Bei
allen nachfolgenden Ergebnissen werden ausschließlich relevante Abweichungen hervorgehoben.
Sämtliche Diagramme sind in Kapitel \ref{chap:appendix} zu finden.

\input{functions/markChart}
\newcommand{\memoryChart}[2][]{%
    % #1 = optional max difficulty (empty -> show all)
    % #2 = CSV file path

    % Load the CSV data
    \pgfplotstableread[col sep=comma]{#2}\loadedMemoryData

    % Store the max difficulty in a new macro
    \ifx\relax#1\relax
        \def\maxDifficulty{0}% Use 0 or some other indicator for no max
        \def\hasMax{0}%
    \else
        \edef\maxDifficulty{#1}% <--- *** Crucial change: Store #1's content ***
        \def\hasMax{1}%
    \fi

    \ifnum\hasMax=1
        % If max is set, \currentXmax is the number
        \edef\currentXmax{\maxDifficulty}
    \else
        % If no max is set, \currentXmax is the pgfplots 'no value' token
        \let\currentXmax\pgfkeysnovalue
    \fi

    \begin{tikzpicture}

        % --- Left axis (ILP memory) ---
        \begin{axis}[
                title={Hauptspeicher (Bytes)},
                name=ilpaxis,
                axis y line*=left,
                xlabel={Difficulty},
                ylabel={ILP Speichernutzung (Bytes)},
                ymode=log,
                % xmax=\currentXmax,
                ymajorgrids=true,
                width=0.95\textwidth,
                height=7cm,
                grid=major,
                legend style={at={(0.5,-0.2)}, anchor=north, legend columns=-1},
                table/x=difficulty,
            ]

            % --- ILP Memory Usage ---
            \addplot[
                mark=*,
                color=blue,
                line width=1pt,
                smooth,
                error bars/.cd,
                y dir=both,
                y explicit,
                /pgfplots/x filter/.code={\ifnum\hasMax=1 \ifdim\pgfmathresult pt>\maxDifficulty pt \def\pgfmathresult{} \fi \fi}
            ] table [
                    x=difficulty,
                    y=ILP,
                    y error expr=\thisrow{ILP stdev},
                ] {\loadedMemoryData};
            \addlegendentry{ILP}

            % --- ILP GPU RAM ---
            \addplot[
                mark=*,
                color=purple,
                line width=1pt,
                smooth,
                error bars/.cd,
                y dir=both,
                y explicit,
                /pgfplots/x filter/.code={\ifnum\hasMax=1 \ifdim\pgfmathresult pt>\maxDifficulty pt \def\pgfmathresult{} \fi \fi}
            ] table [
                    x=difficulty,
                    y={ILP GPU},
                    y error={ILP GPU stdev},
                ] {\loadedMemoryData};
            \addlegendentry{ILP GPU RAM}

            % --- Hill-Climbing v1 ---
            \addplot[
                mark=square*,
                color=red,
                line width=1pt,
                smooth,
                error bars/.cd,
                y dir=both,
                y explicit,
                /pgfplots/x filter/.code={\ifnum\hasMax=1 \ifdim\pgfmathresult pt>\maxDifficulty pt \def\pgfmathresult{} \fi \fi}
            ] table [
                    x=difficulty,
                    y={Hill Climbing v1},
                    y error expr=\thisrow{Hill Climbing v1 stdev},
                ] {\loadedMemoryData};
            \addlegendentry{Hill-Climbing v1}

            % --- Hill-Climbing v3 ---
            \addplot[
                mark=triangle*,
                color=green!60!black,
                line width=1pt,
                smooth,
                error bars/.cd,
                y dir=both,
                y explicit,
                /pgfplots/x filter/.code={\ifnum\hasMax=1 \ifdim\pgfmathresult pt>\maxDifficulty pt \def\pgfmathresult{} \fi \fi}
            ] table [
                    x=difficulty,
                    y={Hill Climbing v3},
                    y error expr=\thisrow{Hill Climbing v3 stdev},
                ] {\loadedMemoryData};
            \addlegendentry{Hill-Climbing v3}

            % --- Offering-Order ---
            \addplot[
                mark=diamond*,
                color=orange,
                line width=1pt,
                dashed,
                smooth,
                error bars/.cd,
                y dir=both,
                y explicit,
                /pgfplots/x filter/.code={\ifnum\hasMax=1 \ifdim\pgfmathresult pt>\maxDifficulty pt \def\pgfmathresult{} \fi \fi}
            ] table [
                    x=difficulty,
                    y={Offering Order},
                    y error expr=\thisrow{Offering Order stdev},
                ] {\loadedMemoryData};
            \addlegendentry{Offering-Order}

        \end{axis}
    \end{tikzpicture}
}

\newcommand{\vramChart}[2][]{%
    % #1 = optional max difficulty (empty -> show all)
    % #2 = CSV file path

    \pgfplotstableread[col sep=comma]{#2}\loadedVramData

    \ifx\relax#1\relax
        \def\vramMaxDifficulty{0}%
        \def\vramHasMax{0}%
    \else
        \edef\vramMaxDifficulty{#1}%
        \def\vramHasMax{1}%
    \fi

    \begin{tikzpicture}
        \begin{axis}[
                title={ILP GPU VRAM (Bytes)},
                axis y line*=left,
                xlabel={Difficulty},
                ylabel={VRAM Nutzung (Bytes)},
                ymajorgrids=true,
                width=0.95\textwidth,
                height=7cm,
                grid=major,
                legend style={at={(0.5,-0.2)}, anchor=north, legend columns=-1},
                table/x=difficulty,
            ]

            \addplot[
                mark=pentagon*,
                color=cyan!70!black,
                line width=1pt,
                dashed,
                smooth,
                error bars/.cd,
                y dir=both,
                y explicit,
                /pgfplots/x filter/.code={\ifnum\vramHasMax=1 \ifdim\pgfmathresult pt>\vramMaxDifficulty pt \def\pgfmathresult{} \fi \fi}
            ] table [
                    x=difficulty,
                    y={ILP GPU VRAM},
                    y error={ILP GPU VRAM stdev},
                ] {\loadedVramData};
            \addlegendentry{ILP GPU VRAM}

        \end{axis}
    \end{tikzpicture}
}

\input{functions/timingChart}


% -----------------------------------------------------------------------
\subsection{Szenario 1: Keine Kurse an Montagen}

Bis Schwierigkeitsstufe 10 findet \textit{Offering-Order} stets den optimalen Stundenplan und
erreicht konstant 100\,\% des ILP-Optimums. \textit{Hill-Climbing v1} und \textit{v3} liefern
in diesem Bereich gleichwertige Ergebnisse: Bei den Stufen 1--3 erzielen beide ebenfalls
100\,\%, sinken jedoch bis Stufe 9 leicht auf 96--97\,\% ab und bleiben damit durchgehend
hinter \textit{Offering-Order}. Ab Schwierigkeitsstufe 10 divergieren die beiden Varianten:
\textit{Hill-Climbing v1} hält sich bei den Stufen 11 und 12 stabil bei rund 94\,\%, während
\textit{Hill-Climbing v3} von 95\,\% bei Stufe 10 auf 90\,\% bei Stufe 11 und weiter auf
ca.\,88\,\% bei Stufe 12 abfällt und bei Stufe 13 keine gültige Lösung mehr findet. Die
Lösungsqualität von \textit{Offering-Order} bricht ab Schwierigkeitsstufe 10 deutlich ein ---
von 100\,\% auf ca.\,90\,\% --- und fällt damit unter jene von \textit{Hill-Climbing v1}.

\vspace{1cm}
\markChart[13]{../reports/csv/score_diff__constraint_4411_no_mondays.csv}

\newpage

Bis Schwierigkeitsstufe 10 steigen die Laufzeiten von \textit{Hill-Climbing v1} und
\textit{v3} gleichmäßig an --- von jeweils rund 0{,}1\,s bei Stufe 1 auf ca.\,0{,}9\,s bei
Stufe 10. ILP weist mit konstant ca.\,0{,}35\,s eine höhere, aber stabile Laufzeit auf, die
selbst bei niedrigen Schwierigkeitsgraden über jener der \textit{Hill-Climbing}-Varianten
liegt. ILP GPU steigt im selben Bereich von 0{,}07\,s auf 0{,}5\,s an. \textit{Offering-Order}
bleibt bis Stufe 10 konstant bei ca.\,0{,}17\,s und ist damit der schnellste Algorithmus.
Auffällig ist, dass ab Schwierigkeitsstufe 10 die Laufzeiten --- mit Ausnahme von
\textit{Hill-Climbing v1} --- zurückgehen: \textit{Hill-Climbing v3} sinkt von 0{,}9\,s auf
0{,}6\,s, ILP von 0{,}35\,s auf 0{,}2\,s, und \textit{Offering-Order} besonders stark von
0{,}18\,s auf 0{,}05\,s.

\vspace{1cm}
\timingChart[13]{../reports/csv/timing_diff__constraint_4411_no_mondays.csv}

\newpage

Bis Schwierigkeitsgrad 10 ist die Hauptspeichernutzung aller Algorithmen konstant, jedoch auf
deutlich unterschiedlichen Niveaus: ILP CPU belegt stabil ca.\,1{,}3\,MB und
\textit{Offering-Order} ca.\,2{,}3\,MB, während \textit{Hill-Climbing v1},
\textit{Hill-Climbing v3} und ILP GPU mit jeweils ca.\,35\,KB deutlich darunter liegen. Ab
Schwierigkeitsgrad 10 verschiebt sich dieses Bild: Die Speichernutzung von ILP GPU steigt auf
ca.\,1{,}3\,MB an und nähert sich damit dem Niveau von ILP CPU, während
\textit{Offering-Order} auf ca.\,35\,KB fällt. Alle übrigen Algorithmen bleiben stabil.

\vspace{1cm}
\memoryChart[13]{../reports/csv/memory_diff__constraint_4411_no_mondays.csv}

\newpage

Die VRAM-Auslastung von ILP GPU beträgt über alle gemessenen Schwierigkeitsstufen konstant
ca.\,14\,Gigabyte.

\vspace{1cm}
\vramChart[13]{../reports/csv/memory_diff__constraint_4411_no_mondays.csv}

\newpage


% -----------------------------------------------------------------------
\subsection{Szenario 2: Keine Kurse an Freitagen}

Die Lösungsqualität von \textit{Hill-Climbing v1} und \textit{v3} entwickelt sich bis
Schwierigkeitsgrad 10 identisch: Ausgehend von 100\,\% bei den Stufen 1--3 sinkt sie
kontinuierlich auf ca.\,94\,\% bei Stufe 10. Ab Schwierigkeitsgrad 11 divergieren die beiden
Varianten --- \textit{Hill-Climbing v3} fällt auf 88\,\%, während \textit{Hill-Climbing v1}
mit 89\,\% leicht besser abschneidet. \textit{Offering-Order} liegt über alle
Schwierigkeitsstufen hinweg unter beiden \textit{Hill-Climbing}-Varianten und sinkt von
95\,\% bei Stufe 2 auf ca.\,86\,\% bei Stufe 13.

\vspace{1cm}
\markChart[14]{../reports/csv/score_diff__constraint_4412_no_fridays.csv}

\newpage

Die Laufzeiten von \textit{Hill-Climbing v1} und \textit{v3} steigen mit zunehmender
Schwierigkeit an: \textit{Hill-Climbing v1} wächst von 0{,}15\,s bei Stufe 1 auf bis zu
0{,}95\,s bei den Stufen 12--14, \textit{Hill-Climbing v3} von 0{,}1\,s auf 0{,}58\,s bei
den Stufen 10--13. \textit{Offering-Order} bleibt durchgehend der schnellste Algorithmus und
hält sich stabil zwischen 0{,}08\,s bei Stufe 1 und 0{,}06\,s bei Stufe 13. ILP CPU und ILP
GPU weisen konstante, aber höhere Laufzeiten auf --- ca.\,0{,}30\,s bzw.\,0{,}67\,s ---, wobei
die Laufzeit von ILP CPU ab Schwierigkeitsgrad 6 leicht zwischen 0{,}33\,s und 0{,}69\,s
schwankt.

\vspace{1cm}
\timingChart[14]{../reports/csv/timing_diff__constraint_4412_no_fridays.csv}

\newpage

Die Hauptspeicherauslastung ist in diesem Szenario über alle Schwierigkeitsstufen hinweg
konstant. ILP CPU belegt mit ca.\,4\,MB den höchsten Speicherbedarf, gefolgt von ILP GPU mit
ca.\,2{,}1\,MB. \textit{Hill-Climbing v1}, \textit{Hill-Climbing v3} und
\textit{Offering-Order} belegen mit ca.\,40\,KB, 33\,KB bzw.\,26\,KB deutlich weniger
Speicher.

\vspace{1cm}
\memoryChart[14]{../reports/csv/memory_diff__constraint_4412_no_fridays.csv}

\newpage


% -----------------------------------------------------------------------
\subsection{Szenario 3: Jeder Tag der Woche bis 9:45 frei}

Die Ergebnisse verlaufen weitgehend analog zu den vorigen Szenarien. Bemerkenswert ist, dass
\textit{Hill-Climbing v3} trotz des reduzierten Suchraums zwischen den Schwierigkeitsstufen 4
und 7 bessere Stundenpläne generiert als \textit{Hill-Climbing v1}: \textit{Hill-Climbing v3}
hält in diesem Bereich konstant 100\,\%, während \textit{Hill-Climbing v1} auf 98--99\,\%
absinkt. Ab Stufe 7 kehrt sich dieses Verhältnis um --- bei Stufe 12 etwa erzielt
\textit{Hill-Climbing v1} 94\,\% gegenüber 90\,\% bei \textit{Hill-Climbing v3}. Auffällig
ist zudem, dass \textit{Offering-Order} bei Schwierigkeitsstufe 14 mit 94\,\% einen besseren
Stundenplan erzeugen konnte als \textit{Hill-Climbing v1} mit 92\,\%.

\vspace{1cm}
\markChart[14]{../reports/csv/score_diff__constraint_4413_everyday_free_until_0945.csv}

\newpage


% -----------------------------------------------------------------------
\subsection{Szenario 4: Keine Kurse an Donnerstagen und Freitagen}

Auch hier verhalten sich die Ergebnisse weitgehend analog zu den vorherigen Szenarien.
Auffällig ist, dass \textit{Hill-Climbing v1} und \textit{v3} über alle Schwierigkeitsstufen
hinweg identische Ergebnisse liefern. \textit{Offering-Order} hält gut mit, liegt bei
einzelnen Schwierigkeitsgraden jedoch leicht darunter: bei Stufe 4 etwa 95\,\% gegenüber
99\,\%, bei Stufe 6 93\,\% gegenüber 95\,\%, bei Stufe 7 94\,\% gegenüber 95\,\% und bei
Stufe 9 91\,\% gegenüber 92{,}5\,\%. Interessant ist, dass \textit{Offering-Order} bei
Schwierigkeitsgrad 11 noch eine gültige Lösung mit 89\,\% findet, während die
\textit{Hill-Climbing}-Varianten an dieser Stelle bereits scheitern.

\vspace{1cm}
\markChart[11]{../reports/csv/score_diff__constraint_4414_free_thursday_friday.csv}

\newpage


% -----------------------------------------------------------------------
\subsection{Szenario 5: Erhöhte Priorität (+90) für Nachmittag, negative Priorität (-90) für Abend}

Bis Schwierigkeitsstufe 5 erreichen alle heuristischen Algorithmen das optimale Ergebnis von
100\,\%. Danach setzt eine Staffelung ein: \textit{Hill-Climbing v1} liefert die längste
Lösungsserie --- bis Stufe 16 hält es ca.\,95\,\%, bevor es bei Stufe 17 auf 88\,\% abfällt.
\textit{Hill-Climbing v3} produziert bis Stufe 9 stabile 99\,\%, sinkt dann jedoch
kontinuierlich und erreicht bei Stufe 15 noch 87\,\%. \textit{Offering-Order} zeigt ein
volatileres Bild: Bei Stufe 6 fällt die \textit{Mark} auf 85\,\% und erreicht bei Stufe 9 mit 70\,\%
ihren niedrigsten Wert, bevor sie sich bei den Stufen 11--15 auf ca.\,88\,\% erholt. Bei Stufe
16 fällt \textit{Offering-Order} erneut auf 72\,\% und findet ab Stufe 17 keine gültige Lösung
mehr.

\vspace{1cm}
\markChart[17]{../reports/csv/score_diff__constraint_4421_priority.csv}

\newpage

Das Laufzeitverhalten ähnelt jenem aus Szenario 2. ILP zeigt bis Schwierigkeitsstufe 13 eine
stabile Laufzeit von ca.\,0{,}55\,s, steigt danach jedoch kontinuierlich an und erreicht bei
Stufe 17 ca.\,2\,s. \textit{Hill-Climbing v1} wächst von unter 1\,s bei den Stufen 1--2 auf
ca.\,3\,s ab Stufe 9 und verbleibt auf diesem Niveau. \textit{Hill-Climbing v3} bleibt bis
Stufe 6 unter 1\,s und stabilisiert sich ab Stufe 10 bei ca.\,1{,}3\,s. ILP GPU ist mit
1{,}1--1{,}2\,s konstant. \textit{Offering-Order} läuft bis Stufe 15 stabil bei ca.\,0{,}17\,s,
schnellt bei Stufe 16 auf 7{,}4\,s hoch und überschreitet bei Stufe 17 das Zeitlimit.

\vspace{1cm}
\timingChart[17]{../reports/csv/timing_diff__constraint_4421_priority.csv}

\newpage


% -----------------------------------------------------------------------
\subsection{Szenario 6: Erhöhte Priorität (+70) für Vormittag, neutrale Priorität (0) für Nachmittag, negative Priorität (-80) für Abend}

Auffällig ist, dass \textit{Hill-Climbing v3} und \textit{Offering-Order} ab
Schwierigkeitsstufe 14 noch gültige Lösungen finden, während \textit{Hill-Climbing v1} an
dieser Stelle bereits keine Ergebnisse mehr produziert. Gleichzeitig bleibt die
Lösungsqualität von \textit{Hill-Climbing v1} in diesem Szenario durchgehend hoch und fällt
nie unter 95\,\%. \textit{Hill-Climbing v3} hingegen bewegt sich zwischen den Stufen 8 und 13
im Bereich von 90--95\,\%, fällt bei Stufe 14 auf 75\,\% und bei Stufe 15 auf 60\,\%.
\textit{Offering-Order} findet bei Stufe 15 noch eine Lösung mit 72\,\%, schwankt aber bei
den Stufen 16 und 17 zwischen 58\,\% und 62\,\%.

\vspace{1cm}
\markChart[17]{../reports/csv/score_diff__constraint_4422_priority.csv}

\newpage

Das Laufzeitverhalten ähnelt jenem aus Szenario 2. \textit{Hill-Climbing v1} steigt von
0{,}5\,s bei Stufe 1 auf ca.\,5\,s ab Stufe 9, wo es sich stabilisiert. ILP GPU bleibt
konstant bei ca.\,1{,}2\,s, während ILP CPU von 0{,}5\,s bei Stufe 1 auf maximal 1{,}3\,s bei
Stufe 16 ansteigt. \textit{Offering-Order} hält bis Schwierigkeitsstufe 14 eine stabile
Laufzeit von ca.\,0{,}26\,s, steigt dann jedoch stark an --- auf 0{,}5\,s bei Stufe 15, 4\,s
bei Stufe 16 und 7\,s bei Stufe 17 --- und überschreitet bei Stufe 18 das Zeitlimit. Dieser
starke Anstieg lässt sich auf häufiges \textit{Backtracking} zurückführen, das ein
\textit{Brute-Force}-ähnliches Suchverhalten nach sich zieht.

\vspace{1cm}
\timingChart[17]{../reports/csv/timing_diff__constraint_4422_priority.csv}

\newpage


% -----------------------------------------------------------------------
\subsection{Szenario 7: Studierende müssen von 12:00 bis 13:00 Zeit für ein Mittagessen haben}

\textit{Hill-Climbing v3} erzeugt zwischen den Schwierigkeitsstufen 6 und 10 trotz
eingeschränktem Suchraum bessere Stundenpläne als \textit{Hill-Climbing v1}: In diesem Bereich
hält \textit{Hill-Climbing v3} konstant 100\,\%, während \textit{Hill-Climbing v1} auf 98\,\%
abfällt. Bei höheren Schwierigkeitsstufen verbessern sich die Ergebnisse von
\textit{Offering-Order} zunehmend: Nach einem Tiefpunkt von 78\,\% bei Stufe 5 steigt die
Qualität kontinuierlich auf 92\,\% bei Stufe 13 und verbleibt auf diesem Niveau.

\vspace{1cm}
\markChart[17]{../reports/csv/score_diff__constraint_4431_lunch.csv}

\newpage

Auch die Laufzeiten verlaufen ähnlich wie in Szenario 2. \textit{Hill-Climbing v1} steigt von
0{,}3\,s bei Stufe 1 auf ca.\,3\,s bei Stufe 9 und verbleibt auf diesem Niveau.
\textit{Hill-Climbing v3} wächst von 0{,}2\,s auf 2\,s bei Stufe 10 und stabilisiert sich
dort. ILP CPU und ILP GPU sind mit ca.\,0{,}22\,s bzw.\,0{,}38\,s konstant, wobei ILP GPU
leicht höhere Laufzeiten aufweist. Bemerkenswert ist, dass die Laufzeit von
\textit{Offering-Order} mit steigender Schwierigkeitsstufe sogar leicht abnimmt: von 0{,}17\,s
bei den Stufen 2--8 auf 0{,}13\,s bei den Stufen 9--14 und schließlich 0{,}1\,s bei den
Stufen 15--17.

\vspace{1cm}
\timingChart[17]{../reports/csv/timing_diff__constraint_4431_lunch.csv}

\newpage


% -----------------------------------------------------------------------
\subsection{Szenario 8: Studierende müssen mehr als einen Kurs pro Tag besuchen}

Die Qualität der erstellten Stundenpläne schwankt in diesem Szenario stark.
\textit{Hill-Climbing v1} und \textit{v3} liefern bis Schwierigkeitsgrad 6 gleichwertige
Ergebnisse, die zwischen 100\,\% bei Stufe 1 und 67\,\% bei Stufe 6 schwanken. Bei
Schwierigkeitsgrad 7 liegt \textit{Hill-Climbing v1} mit 89\,\% knapp vor
\textit{Hill-Climbing v3} mit 88\,\%. \textit{Offering-Order} zeigt das volatilste Verhalten:
Die Qualität startet bei 70\,\% bei Stufe 1 und schwankt in der Folge zwischen 69\,\% bei
Stufe 8 und 98\,\% bei Stufe 7.

\vspace{1cm}
\markChart[9]{../reports/csv/score_diff__constraint_4432_more_than_one_course_a_day.csv}

\newpage

Die Laufzeiten von ILP und ILP GPU schwanken in diesem Szenario erheblich: ILP bewegt sich
zwischen 0{,}1\,s bei Stufe 5 und 5{,}7\,s bei Stufe 8, ILP GPU zwischen 1{,}2\,s bei Stufe
1 und 6{,}7\,s bei Stufe 5. \textit{Hill-Climbing v1} und \textit{v3} steigen moderat an ---
von 0{,}04\,s bzw.\,0{,}03\,s bei Stufe 1 auf 0{,}14\,s bzw.\,0{,}09\,s bei Stufe 7.
\textit{Offering-Order} bleibt bis Schwierigkeitsgrad 8 konstant bei 0{,}13\,s und steigt
erst bei Schwierigkeitsgrad 9 auf 0{,}6\,s an.

\vspace{1cm}
\timingChart[9]{../reports/csv/timing_diff__constraint_4432_more_than_one_course_a_day.csv}

\newpage

Der Hauptspeicherbedarf von ILP schwankt in diesem Szenario stark: Im Regelfall belegt ILP
ca.\,8{,}4\,MB, fällt jedoch bei Schwierigkeitsstufe 5 auf 38\,KB und bei Stufe 7 auf
700\,KB ab. ILP GPU bewegt sich zwischen 4{,}2\,MB und 8{,}4\,MB.
\textit{Hill-Climbing v1}, \textit{Hill-Climbing v3} und \textit{Offering-Order} sind mit
15--23\,KB, 23--33\,KB bzw.\,29--41\,KB konstant auf niedrigem Niveau.

\vspace{1cm}
\memoryChart[9]{../reports/csv/memory_diff__constraint_4432_more_than_one_course_a_day.csv}

\newpage


% -----------------------------------------------------------------------
\subsection{Szenario 9: Studierende dürfen nicht mehr als drei Kurse pro Tag besuchen}

Die Lösungsqualität von \textit{Hill-Climbing v1} und \textit{v3} ist bis Schwierigkeitsstufe
13 gleichwertig: Ausgehend von 100\,\% bei Stufe 1 sinkt sie kontinuierlich auf 96\,\% bei
Stufe 13 ab. Danach erzielt \textit{Hill-Climbing v1} leicht bessere Ergebnisse --- bei Stufe
16 etwa 95\,\% gegenüber 94\,\% bei \textit{Hill-Climbing v3}. \textit{Offering-Order} liegt
über alle Schwierigkeitsstufen hinweg unter den \textit{Hill-Climbing}-Varianten und erreicht
bei den Stufen 8 und 9 mit 88\,\% seinen niedrigsten Wert, der damit als einziger unter die
90\,\%-Marke fällt. Generell ist die Lösungsqualität in diesem Szenario sehr gut.

\vspace{1cm}
\markChart[17]{../reports/csv/score_diff__constraint_4433_max_3_courses.csv}

\newpage

Die Laufzeit von ILP CPU steigt kontinuierlich von 0{,}3\,s bei Stufe 1 auf 1{,}8\,s bei
Stufe 17, sodass ab Stufe 14 ILP GPU mit konstant ca.\,1\,s die kürzeren Laufzeiten aufweist.
\textit{Hill-Climbing v1} und \textit{v3} steigen ebenfalls mit der Schwierigkeit an --- von
0{,}2\,s bzw.\,0{,}1\,s bei Stufe 1 auf 2\,s bzw.\,1\,s bei Stufe 17. \textit{Offering-Order}
bleibt mit ca.\,0{,}1\,s konstant auf dem niedrigsten Niveau aller Algorithmen.

\vspace{1cm}
\timingChart[17]{../reports/csv/timing_diff__constraint_4433_max_3_courses.csv}

\newpage


% -----------------------------------------------------------------------
\subsection{Szenario 10: Alle Kurse können frei von Constraints gewählt werden}

Die Qualität der Stundenpläne ist mit jener aus Szenario 9 vergleichbar. Die Lösungsqualität von \textit{Hill-Climbing v1} und \textit{v3} ist bis Schwierigkeitsstufe
13 gleichwertig: Ausgehend von 100\,\% bei Stufe 1 sinkt sie kontinuierlich auf 96\,\% bei
Stufe 13 ab. Danach erzielt \textit{Hill-Climbing v1} leicht bessere Ergebnisse --- bei Stufe
16 etwa 95\,\% gegenüber 94\,\% bei \textit{Hill-Climbing v3}. \textit{Offering-Order} liegt
über alle Schwierigkeitsstufen hinweg unter den \textit{Hill-Climbing}-Varianten und erreicht
bei den Stufen 8 und 9 mit 88\,\% seinen niedrigsten Wert, der damit als einziger unter die
90\,\%-Marke fällt. Generell ist die Lösungsqualität in diesem Szenario sehr gut.

\markChart[17]{../reports/csv/score_diff__constraint_4441_no_constraints.csv}
\vspace{1cm}

\newpage

Die Laufzeit von ILP CPU steigt kontinuierlich von 0{,}3\,s bei Stufe 1 auf 1{,}8\,s bei
Stufe 17, sodass ab Stufe 14 ILP GPU mit konstant ca.\,1\,s die kürzeren Laufzeiten aufweist.
\textit{Hill-Climbing v1} und \textit{v3} steigen ebenfalls mit der Schwierigkeit an --- von
0{,}2\,s bzw.\,0{,}1\,s bei Stufe 1 auf 2\,s bzw.\,1\,s bei Stufe 17. \textit{Offering-Order}
bleibt mit ca.\,0{,}1\,s konstant auf dem niedrigsten Niveau aller Algorithmen.

\vspace{1cm}
\timingChart[17]{../reports/csv/timing_diff__constraint_4441_no_constraints.csv}

\newpage


% -----------------------------------------------------------------------
\subsection{Szenario 11: 25\% -- 75\% aller Kurse können nicht belegt werden}

Beide \textit{Hill-Climbing}-Varianten erzielen über alle Schwierigkeitsstufen hinweg
identische Ergebnisse: Ausgehend von 100\,\% bei Stufe 1 sinkt die Lösungsqualität
kontinuierlich auf 94\,\% bei Stufe 15. \textit{Offering-Order} liegt ab Schwierigkeitsgrad 2
durchgehend darunter und bewegt sich zwischen den Stufen 4 und 15 im Bereich von 85--90\,\%.

\vspace{1cm}
\markChart[15]{../reports/csv/score_diff__constraint_4442_25_pct_blocked_1.csv}

\newpage

Die Laufzeiten von \textit{Hill-Climbing v1} und \textit{v3} steigen kontinuierlich mit der
Schwierigkeit: \textit{Hill-Climbing v1} wächst von 0{,}13\,s bei Stufe 1 auf 1{,}2\,s bei
Stufe 15, \textit{Hill-Climbing v3} von 0{,}09\,s auf 0{,}7\,s. \textit{Offering-Order} bleibt
bis Schwierigkeitsgrad 14 konstant bei ca.\,0{,}07\,s und steigt bei Stufe 15 deutlich auf
0{,}4\,s an --- womit es über alle Stufen der schnellste Algorithmus ist. ILP GPU ist mit
ca.\,0{,}7\,s konstant. ILP CPU zeigt bis Schwierigkeitsgrad 6 eine stabile Laufzeit von
ca.\,0{,}36\,s und steigt danach leicht auf 0{,}5--0{,}6\,s an.

\vspace{1cm}
\timingChart[15]{../reports/csv/timing_diff__constraint_4442_25_pct_blocked_1.csv}

\newpage


% -----------------------------------------------------------------------
\subsection{Szenario 12: 1 -- 7 Kurse sind vorgegeben}

\textit{Offering-Order} zeigt in diesem Szenario besonders schwache Leistungen: Bei Stufe 1
erreicht der Algorithmus noch 100\,\%, fällt bei Stufe 2 jedoch abrupt auf 30\,\% und hält
sich bei den Stufen 3 und 4 bei ca.\,27\,\%, bevor er ab Schwierigkeitsstufe 5 das Zeitlimit
überschreitet und keine Ergebnisse mehr produziert. \textit{Hill-Climbing v1} und \textit{v3}
liefern bis Schwierigkeitsstufe 12 gleichwertige Ergebnisse im Bereich von 100\,\% bis
ca.\,95\,\%; danach erzeugt \textit{Hill-Climbing v1} leicht bessere Stundenpläne --- bei
Stufe 16 etwa 93\,\% gegenüber 91\,\% bei \textit{Hill-Climbing v3}.

\vspace{1cm}
\markChart[17]{../reports/csv/score_diff__constraint_4443_1_scheduled_1.csv}

\newpage

Die Laufzeiten beider \textit{Hill-Climbing}-Varianten steigen kontinuierlich mit zunehmender
Schwierigkeit: \textit{Hill-Climbing v1} wächst von 0{,}18\,s bei Stufe 1 auf 2{,}7\,s bei
Stufe 17, \textit{Hill-Climbing v3} von 0{,}1\,s auf 1{,}3\,s bei Stufe 16.
\textit{Offering-Order} weist mit konstant 0{,}13\,s stabil geringe Laufzeiten auf, läuft
jedoch bereits ab Schwierigkeitsstufe 4 ins Zeitlimit.

\vspace{1cm}
\timingChart[17]{../reports/csv/timing_diff__constraint_4443_1_scheduled_1.csv}

\newpage


% -----------------------------------------------------------------------
\subsection{Szenario 13: Keine Kurse zwischen 6:00 Uhr morgens und 13:00 Uhr}

\textit{Hill-Climbing v3} erzielt in diesem Szenario ab Schwierigkeitsstufe 3 durchgehend
bessere Ergebnisse als \textit{Hill-Climbing v1}: Bei den Stufen 1--3 erzielen beide noch
100\,\%; bei den Stufen 4 und 5 hält \textit{Hill-Climbing v3} 100\,\%, während
\textit{Hill-Climbing v1} bereits zurückfällt. Dieser Vorsprung wächst mit der Schwierigkeit
--- bei Stufe 10 etwa erreicht \textit{Hill-Climbing v3} 92\,\%, \textit{Hill-Climbing v1}
nur 83\,\%. Bemerkenswert ist, dass \textit{Hill-Climbing v3} und \textit{Offering-Order} bei
Schwierigkeitsstufe 11 noch gültige Lösungen finden, obwohl \textit{Hill-Climbing v1} hier
bereits scheitert --- und das, obwohl der Suchraum bei \textit{Hill-Climbing v3} eingeschränkt
ist. \textit{Offering-Order} bewegt sich bei den Stufen 1--3 bei 100\,\% und schwankt danach
zwischen 82\,\% (Stufe 10) und 92\,\% (Stufe 7).

\vspace{1cm}
\markChart[11]{../reports/csv/score_diff__constraint_4444_no_course_between_6_and_13.csv}

\newpage

Die Laufzeiten der \textit{Hill-Climbing}-Varianten steigen gleichmäßig mit zunehmender
Schwierigkeit: \textit{Hill-Climbing v1} wächst von 0{,}13\,s bei Stufe 1 auf 0{,}84\,s bei
Stufe 10, \textit{Hill-Climbing v3} von 0{,}11\,s auf 0{,}7\,s. \textit{Offering-Order} ist
die schnellste Heuristik und hält bei den Stufen 1--6 ca.\,0{,}07\,s, sinkt bis Stufe 10
sogar leicht auf 0{,}04\,s, bevor es bei Stufe 11 mit 2{,}4\,s einen deutlichen Anstieg
verzeichnet.

\vspace{1cm}
\timingChart[11]{../reports/csv/timing_diff__constraint_4444_no_course_between_6_and_13.csv}

\newpage


% -----------------------------------------------------------------------
\subsection{Szenario 14: Keine Kurse am Montag, Dienstag und Mittwoch}

\textit{Offering-Order} erzielt in diesem Szenario bei den Schwierigkeitsstufen 7 und 8 noch
gültige Lösungen, während \textit{Hill-Climbing v1} und \textit{v3} an dieser Stelle bereits
scheitern. \textit{Hill-Climbing v1} erreicht bei den Stufen 1--4 noch 100\,\% und sinkt bis
Stufe 6 auf 94\,\%, findet danach jedoch keine Lösung mehr. \textit{Hill-Climbing v3} fällt
nach 100\,\% bei den Stufen 1--3 bereits bei Stufe 4 auf 90\,\% und bei Stufe 5 auf 86\,\% ab.
\textit{Offering-Order} hält bei den Stufen 1--4 ebenfalls 100\,\%, erreicht bei den Stufen
6 und 7 seinen niedrigsten Wert von 94\,\% und steigt bei Stufe 8 wieder auf 96\,\%.

\vspace{1cm}
\markChart[8]{../reports/csv/score_diff__constraint_4445_no_course_monday_tuesday_wednesday.csv}

\newpage

Die Laufzeit von ILP CPU ist mit konstant ca.\,0{,}03\,s durchgehend niedriger als jene von
ILP GPU mit ca.\,0{,}18\,s. \textit{Hill-Climbing v1} steigt von 0{,}03\,s bei Stufe 1 auf
0{,}07\,s bei Stufe 4 und verbleibt dort bis Stufe 6. \textit{Hill-Climbing v3} wächst von
0{,}03\,s auf 0{,}04\,s bei Stufe 2 und stabilisiert sich bis Stufe 5 auf diesem Niveau.
\textit{Offering-Order} ist die schnellste Heuristik und läuft konstant bei ca.\,0{,}01\,s,
verzeichnet bei Schwierigkeitsstufe 8 jedoch einen Anstieg auf 0{,}18\,s.

\vspace{1cm}
\timingChart[8]{../reports/csv/timing_diff__constraint_4445_no_course_monday_tuesday_wednesday.csv}

\newpage